\chapter{Build: A Pricing Library}\label{ch:dv:build-a-pricing-library}

A risk manager asks for the book's \omterm{def:dv:greeks-and-the-hedging-pnl:greeks}{vega} to one node of the \omterm{def:dv:implied-volatility-and-its-surface:surface}{volatility surface}: the one-year, 90\% strike. The book holds a vanilla call, an
American put, a knock-out, a \omterm{def:dv:variance-swaps-and-volatility-derivatives:vs}{variance swap} and an \omterm{def:dv:autocallables:autocallable}{autocallable}. Each was priced by the chapter that introduced it, with its own inputs: a volatility
here, a rate there, a grid of paths somewhere else. None of them reads the surface as an object that can be bumped at one node. The answer needs every
pricer to take the same market snapshot, apply the same bump, and reprice. With a Monte Carlo pricer in the book, it also needs the bumped and unbumped
runs to use the same random numbers, or the answer is noise. This chapter builds the library that makes the question a single call: instruments,
models and engines behind one interface, market data as an immutable snapshot, \omterm{def:dv:greeks-and-the-hedging-pnl:greeks}{Greeks} and scenarios by \omterm{def:dv:greeks-and-the-hedging-pnl:bump}{bump-and-reprice}, and the tests that let other
books trust it. On this chapter's book the node's \omterm{def:dv:greeks-and-the-hedging-pnl:greeks}{vega} is 2\,072 per volatility point. Its Monte Carlo part has a standard deviation of 51 with common
random numbers and 217 without.

\section{Instruments, models and engines}

A pricing library separates three things that a script mixes. The \emph{instrument} is the contract: its terms and nothing else. The \emph{model} is a
set of assumptions about the dynamics, with its parameters. The third is the numerical method.

\begin{definition}[Pricing engine]\label{def:dv:build-a-pricing-library:engine}
A \emph{pricing engine}\index{pricing engine} is the component that computes an instrument's value under a model from a market snapshot by a given
numerical method (closed form, tree, finite differences, Monte Carlo, Fourier); several engines may support the same instrument and model, and a
registry chooses the default.
\end{definition}

The separation pays three times. A new instrument needs only its terms and one engine. A new method, such as a GPU Monte Carlo, prices every instrument
it supports without touching them. And two engines on the same instrument check each other. QuantLib, the open-source library, is built this way: its
abstract \texttt{Instrument} class takes a \omterm{def:dv:build-a-pricing-library:engine}{pricing engine} through \texttt{setPricingEngine} and returns its value through \texttt{NPV}. The Open Source
Risk Engine extends it with further models, instruments and engines and reads trades and market data as XML.

\begin{center}\small
\begin{tabular}{lrrrr}
\toprule
one-year option (skewed surface) & Analytic & Tree & PDE & Monte Carlo\\
\midrule
call, strike 100 & 8.5437 & 8.5437 & 8.5413 & 8.5415\\
American put, strike 95 & --- & 5.0731 & 5.0702 & ---\\
\bottomrule
\end{tabular}
\end{center}

The registry returns every engine that supports a pair, and the tests price across them: the analytic and tree engines agree to $2\times10^{-6}$ on the
call. The grid's 200 nodes leave 0.0024, and one Monte Carlo run's error is of the order of its standard error. For the whole library the engines are the
closed forms of chapters 3, 15 and 16, the trees of chapter 2, the finite differences of chapter 22, the Monte Carlo of chapter 23, the Fourier pricer of
chapter 24 and the convertible grid of chapter 21 (\cref{fig:dv:build-a-pricing-library:arch}).

\begin{omfigure}
\begin{tikzpicture}[font=\small,box/.style={draw,rounded corners,align=center,minimum height=0.9cm,inner sep=4pt},>=stealth]
\node[box,fill=omDef!15,text width=3.2cm] (inst) at (0,2.2) {Instrument\\{\footnotesize terms only, JSON}};
\node[box,fill=omThm!15,text width=3.2cm] (model) at (0,0) {Model\\{\footnotesize Black--Scholes, Heston, credit}};
\node[box,fill=omProp!15,text width=3.4cm] (md) at (0,-2.2) {Market-data snapshot\\{\footnotesize immutable; \texttt{apply(bump)}}};
\node[box,text width=2.6cm] (reg) at (4.4,1.1) {Registry\\{\footnotesize (instrument, model) $\to$ engines}};
\node[box,fill=gray!10,text width=3.3cm] (eng) at (4.4,-1.1) {Engine\\{\footnotesize closed form, tree, PDE, Monte Carlo, Fourier}};
\node[box,text width=3.3cm] (calls) at (9.0,0) {\texttt{price}, \texttt{greeks}, \texttt{sensitivities}, \texttt{bucketed\_vega}, \texttt{reprice}, \texttt{price\_batch}};
\draw[->] (inst) -- (reg);
\draw[->] (model) -- (reg);
\draw[->] (reg) -- (eng);
\draw[->] (md) -- (eng);
\draw[->] (eng) -- (calls);
\draw[->,dashed] (calls.south) |- node[pos=0.75,below,font=\footnotesize]{bumped snapshots} (md.east);
\end{tikzpicture}

\omcaption{The library's structure: an instrument and a model select an engine through the registry; every engine reads the same immutable snapshot; the
top-level calls produce \omterm{def:dv:greeks-and-the-hedging-pnl:greeks}{Greeks}, sensitivities and scenarios by repricing on bumped copies of the snapshot.}\label{fig:dv:build-a-pricing-library:arch}
\end{omfigure}

\section{Market data as an immutable snapshot}

\begin{definition}[Market-data snapshot]\label{def:dv:build-a-pricing-library:snapshot}
A \emph{market-data snapshot}\index{market-data snapshot} is the complete, immutable set of market inputs of a valuation at one time (spots, curves,
dividends, \omterm{def:dv:implied-volatility-and-its-surface:surface}{volatility surfaces}, correlations, fixings, exchange rates), addressable by risk-factor identifiers; a bump or a change of date returns a new
snapshot and never alters the old one.
\end{definition}

QuantLib's instruments cache their last result and observe their inputs, recalculating when an input changes. That design suits a trading screen that
reprices as quotes tick. A risk engine reprices the same book under thousands of shifted markets, often in parallel, and needs the opposite: pricing as a
pure function of an instrument and a snapshot, with no state to invalidate. The library's snapshot has a method \texttt{apply(bump)} that returns a new
snapshot. Bumps are keyed by the risk-factor identifiers that Book 6's risk engine shares: \texttt{SPOT:ABC}, \texttt{VOL:ABC} (parallel),
\texttt{VOL:ABC:2027-09-24:90} (one node), \texttt{CURVE:USD}, \texttt{CORR:ABC|XYZ}, \texttt{TIME}. A surface that has nodes exposes them as risk factors,
and one that does not is first sampled onto a grid.

Immutability is what makes the rest simple. A Greek is two prices on two snapshots, a scenario is one price on one snapshot, and a batch is a matrix of
snapshots by trades. Every engine, old or new, gets all three for free, because none of them knows that bumps exist.

\section{Greeks, bumps and scenarios}

\omterm{def:dv:greeks-and-the-hedging-pnl:bump}{Bump-and-reprice} is the universal Greek: slower than an engine's own derivatives, but available for every engine and consistent across them. Its enemy
is Monte Carlo noise. Two independent runs, each with a standard error of a few tenths, differ by that much, and a bump of one volatility point divides
the difference by 0.01. The library's binding rule removes it. Every Monte Carlo price of an instrument uses the same seed, derived from the
instrument's identifier by a checksum, and the same time grid. The bumped and base runs share their random numbers, and the noise of the difference is
only the noise of the sensitivity itself.

\begin{example}[The vega nobody could compute]\label{ex:dv:build-a-pricing-library:node}
The book, on a surface with 20\% at the money and a skew: long 1\,000 one-year 100 calls, short 500 one-year 95 American puts, long 800 one-year 100 calls
knocked out at 80 (daily monitoring), short a one-year \omterm{def:dv:variance-swaps-and-volatility-derivatives:vs}{variance swap} struck at 21 on a \omterm{def:dv:variance-swaps-and-volatility-derivatives:veganot}{vega notional} of 50\,000, and short 1\,000\,000 of a three-year
\omterm{def:dv:autocallables:autocallable}{autocallable}. Its value is $-928\,523$. The \omterm{def:dv:greeks-and-the-hedging-pnl:greeks}{vega} to the (one year, 90) node, per volatility point, is $-88$ from the put, whose strike interpolates halfway
to the node, $-90$ from the \omterm{def:dv:variance-swaps-and-volatility-derivatives:vs}{variance swap}'s strip and $+2\,249$ from the \omterm{def:dv:autocallables:autocallable}{autocallable}, which the Monte Carlo engine prices on the total variance at the
midpoint of its knock-in (80) and trigger (100): 2\,072 in all. The call and the
knock-out, struck at 100, have none. Over twenty seeds the \omterm{def:dv:autocallables:autocallable}{autocallable}'s contribution has a standard deviation of 51 with common random numbers and 217
without, 4.3 times more (\cref{fig:dv:build-a-pricing-library:crn}).
\end{example}

\begin{omfigure}
\begin{tikzpicture}
\begin{axis}[width=10.5cm,height=5cm,xlabel={seed},ylabel={(1y, 90) vega (thousands a point)},
  tick label style={font=\small},label style={font=\small},
  xmin=0,xmax=21,ymin=1.8,ymax=2.7,grid=major,grid style={gray!25},
  legend columns=2,legend cell align=left,legend style={font=\footnotesize,at={(0.5,-0.3)},anchor=north,draw=none,fill=none,/tikz/every even column/.append style={column sep=8pt}},
  table/col sep=comma]
\addplot[omThm,only marks,mark=*,mark size=1.8pt] table[x=seed,y=with]{figdata/derivatives/28-build-a-pricing-library/crn.csv};
\addplot[omProp,only marks,mark=triangle*,mark size=2pt] table[x=seed,y=without]{figdata/derivatives/28-build-a-pricing-library/crn.csv};
\legend{common random numbers,fresh numbers for each snapshot}
\end{axis}
\end{tikzpicture}

\omcaption{The \omterm{def:dv:autocallables:autocallable}{autocallable}'s \omterm{def:dv:greeks-and-the-hedging-pnl:greeks}{vega} to the one-year 90 node by \omterm{def:dv:greeks-and-the-hedging-pnl:bump}{bump-and-reprice}, from twenty seeds: with the bumped and base runs sharing their random
numbers, and with every snapshot drawing afresh. Data: the tutorial.}\label{fig:dv:build-a-pricing-library:crn}
\end{omfigure}

The node grid gives the whole profile (\cref{fig:dv:build-a-pricing-library:vega}). Summed over its 35 nodes the book's \omterm{def:dv:greeks-and-the-hedging-pnl:greeks}{vega} is 6\,747 per point, against
6\,706 for a parallel bump. The difference comes from interpolation, which spreads a node's bump to its neighbours unevenly.

\begin{omfigure}
\begin{tikzpicture}
\begin{groupplot}[group style={group size=2 by 1,horizontal sep=1.6cm},
  width=6.6cm,height=5cm,
  tick label style={font=\small},label style={font=\small},grid=major,grid style={gray!25},table/col sep=comma]
\nextgroupplot[ybar,bar width=10pt,xlabel={expiry (months)},ylabel={vega (thousands a point)},xtick=data,ymin=-0.5,ymax=3.5,title={\small by expiry}]
\addplot[omDef,fill=omDef!40] table[x=months,y=vega]{figdata/derivatives/28-build-a-pricing-library/vega_expiry.csv};
\nextgroupplot[ybar,bar width=10pt,xlabel={strike},xtick=data,ymin=-0.5,ymax=3.5,title={\small one-year expiry, by strike}]
\addplot[omThm,fill=omThm!40] table[x=strike,y=vega]{figdata/derivatives/28-build-a-pricing-library/vega_1y.csv};
\end{groupplot}
\end{tikzpicture}

\omcaption{The book's bucketed \omterm{def:dv:greeks-and-the-hedging-pnl:greeks}{vega} by \omterm{def:dv:greeks-and-the-hedging-pnl:bump}{bump-and-reprice} at every node of the surface: by expiry (left) and across the one-year slice (right), where the 90
node carries the \omterm{def:dv:autocallables:autocallable}{autocallable}'s exposure. Data: the tutorial.}\label{fig:dv:build-a-pricing-library:vega}
\end{omfigure}

Scenarios are bumps in combination. A scenario is a name and a tuple of bumps, serialisable, so that Book 6's scenario generators can build them and the
library can apply them. The book's ten-scenario grid (\cref{fig:dv:build-a-pricing-library:scen}) shows it long volatility and long a crash, through the
\omterm{def:dv:autocallables:autocallable}{autocallable} it sold, and losing on a rally.

\begin{omfigure}
\begin{tikzpicture}
\begin{axis}[width=10.5cm,height=6cm,xbar,bar width=7pt,xlabel={P\&L of the book (thousands)},
  tick label style={font=\small},label style={font=\small},
  xmin=-50,xmax=175,ymin=-0.6,ymax=9.6,ytick={0,1,2,3,4,5,6,7,8,9},
  yticklabels={spot $-20\%$,spot $-10\%$,spot $-5\%$,spot $+5\%$,spot $+10\%$,vol $-5$,vol $+5$,rates $+100$\,bp,crash: $-20\%$ and vol $+10$,one month later},
  y dir=reverse,grid=major,grid style={gray!25},table/col sep=comma]
\addplot[omProp,fill=omProp!40] table[x=pnl,y=i]{figdata/derivatives/28-build-a-pricing-library/scenarios.csv};
\end{axis}
\end{tikzpicture}

\omcaption{The book repriced under ten scenarios in one call. Data: the tutorial.}\label{fig:dv:build-a-pricing-library:scen}
\end{omfigure}

\section{Testing a pricer}

A pricer is trusted because of its tests, and a library is only as good as the tests of its weakest engine.

\begin{definition}[Reference pricer]\label{def:dv:build-a-pricing-library:reference}
A \emph{reference pricer}\index{reference pricer} is an independent, slow and accurate implementation of a valuation, kept to test production engines against:
a closed form where one exists, otherwise a converged grid or a large simulation, with its own accuracy known.
\end{definition}

The library's tests fall into five kinds. \textbf{Reference values}: each engine against a closed form or a \omterm{def:dv:build-a-pricing-library:reference}{reference pricer} (Black's formula, the Gauss--Legendre
Lewis integral of chapter 24, the tree at 2\,001 steps). \textbf{Cross-engine}: every engine that supports an instrument against the others, within their known
errors. \textbf{Identities}: knock-in plus knock-out equals the vanilla, a digital call plus a digital put pays the discount factor, a one-asset basket equals the
vanilla, a TARF with no leverage and no target equals a strip of puts, higher correlation cheapens a worst-of put. \textbf{Properties}: snapshots are unchanged by
bumps, bucketed \omterm{def:dv:greeks-and-the-hedging-pnl:greeks}{vega} sums to the parallel \omterm{def:dv:greeks-and-the-hedging-pnl:greeks}{vega}, the seed is the same in every process, JSON round-trips every instrument. \textbf{Twins}: the C++20 and Rust cores
reproduce the Python engines to $10^{-9}$.


\section{Tutorial: one call for the whole book}\label{tut:dv:build-a-pricing-library:tut}

\begin{tutorial}
\textbf{Goal.} Wire the book's components into one library, price a small book through one call, compute bucketed \omterm{def:dv:greeks-and-the-hedging-pnl:greeks}{vega} with common random numbers, and run a
ten-scenario grid. \textbf{End state:} the four figures, the engine table and the numbers of the weekend problem.
\begin{tutsteps}
\item \textbf{The registry}: engines by instrument and model type, the first that supports being the default:
\omcode{code/firm/pricing/firm_pricing.py}{506}{521}{The engine registry.}\label{lst:dv:build-a-pricing-library:registry}
\item \textbf{Common random numbers}: the seed from the identifier, and paths drifted to the snapshot's forwards:
\omcode{code/firm/pricing/firm_pricing.py}{1002}{1030}{Seeds and paths of the Monte Carlo engine.}\label{lst:dv:build-a-pricing-library:mc}
\item \textbf{The C++20 twin} of the finite-difference engine:
\omcode{code/firm/pricing/cpp/pricing.hpp}{246}{258}{The PDE engine in C++20.}\label{lst:dv:build-a-pricing-library:cpp}
\item \textbf{Run} \texttt{dv\_library.book\_prices()}, \texttt{engine\_table()}, \texttt{node\_vega()}, \texttt{crn\_noise()}, \texttt{bucket\_profile()},
\texttt{scenario\_grid()} and \texttt{fig\_library.py}.
\end{tutsteps}
\textbf{What to change next.} Price the \omterm{def:dv:autocallables:autocallable}{autocallable} under Heston by adding a Heston Monte Carlo engine; register a second Monte Carlo engine with a million paths as
a \omterm{def:dv:build-a-pricing-library:reference}{reference pricer}; time \texttt{price\_batch} on a hundred snapshots.
\end{tutorial}

\section{Build: the library}\label{bld:dv:build-a-pricing-library:pricing}

\begin{build}
\bfield{Purpose} The miniature firm's pricing library: every instrument of Book 5 behind one interface, the one Book 6's risk engine consumes.
\bfield{Interface} As frozen for Book 6 (\texttt{code/firm/INTERFACES.md}, section 1): \texttt{MarketData} with \texttt{apply}, \texttt{rolled},
\texttt{risk\_factors}; \texttt{Instrument} and \texttt{instrument\_from\_dict}; \texttt{Model}, \texttt{Engine}, \texttt{register}; \texttt{price},
\texttt{greeks}, \texttt{sensitivities}, \texttt{bucketed\_vega}, \texttt{reprice}, \texttt{price\_batch}; \texttt{seed\_for}. Added here: instruments
\texttt{DigitalOption}, \texttt{BarrierOption}, \texttt{AsianOption}, \texttt{ForwardStart}, \texttt{Cliquet}, \texttt{VarianceSwap}, \texttt{BasketOption},
\texttt{WorstOf}, \texttt{Autocallable}, \texttt{TARF}, \texttt{ConvertibleBond}; models \texttt{HestonModel}, \texttt{CreditBlackScholes}; engines
\texttt{ClosedFormEngine}, \texttt{MonteCarloEngine(n\_paths, seed\_offset, crn)}, \texttt{PDEEngine} (with \texttt{fd\_vanilla}), \texttt{FourierEngine},
\texttt{ConvertibleEngine}; \texttt{engines\_for(inst, model)}. The C++20 twin \texttt{cpp/pricing.hpp} and the Rust crate \texttt{rust/} hold the core types
with the analytic and finite-difference engines and the desk \omterm{def:dv:greeks-and-the-hedging-pnl:greeks}{Greeks}.
\bfield{Rules} Pricing is a pure function of instrument, model and snapshot; Monte Carlo seeds come from the instrument's identifier; every engine reads its
market only from the snapshot; additions never rename or remove.
\bfield{Acceptance tests} \texttt{code/firm/pricing/tests/} (the core's fourteen and this chapter's five groups above); \texttt{cpp/pricing\_test.cpp} and
\texttt{rust/src/lib.rs} against the Python values.
\bfield{Stretch} Engines' own \omterm{def:dv:greeks-and-the-hedging-pnl:greeks}{Greeks} (pathwise and adjoint, chapter 23) behind the \texttt{greeks} override; local and stochastic volatility Monte Carlo; seasoned
path-dependent trades from fixings; a service wrapper that prices batches in parallel processes.
\end{build}

\begin{dated}{2026-09}{Trade-representation standards}\label{dat:dv:build-a-pricing-library:standards}
FpML, the XML standard for the electronic dealing and processing of derivatives, is at version 5.13 as a Recommendation, with 5.14 a Last Call Working Draft of
24 August 2026. The Common Domain Model, a model of financial products, trades and their lifecycle events published as code in several languages, has been open
source at FINOS since February 2023, in partnership with ISDA, ICMA and ISLA. A production library maps its instruments to these representations at its edge,
not inside its engines.
\end{dated}

\begin{omsources}
\item QuantLib reference documentation, class \texttt{QuantLib::Instrument} (version 1.43), quantlib.org.
\item Open Source Risk Engine, repository README, github.com/OpenSourceRisk/Engine.
\item FINOS, Common Domain Model, github.com/finos/common-domain-model.
\item FpML, fpml.org.
\end{omsources}

\section{Exercises}

\begin{exercise}[$\star$]\label{exo:dv:build-a-pricing-library:1}
Why must the Monte Carlo seed come from a checksum of the identifier and not from Python's \texttt{hash()}?
\end{exercise}

\begin{exercise}[$\star$]\label{exo:dv:build-a-pricing-library:2}
Why do the call and the knock-out, both struck at 100, have no \omterm{def:dv:greeks-and-the-hedging-pnl:greeks}{vega} to the (one year, 90) node?
\end{exercise}

\begin{exercise}[$\star$]\label{exo:dv:build-a-pricing-library:3}
Name the identity each of these tests checks: knock-in plus knock-out; a digital call plus a digital put; a TARF with no leverage and no target.
\end{exercise}

\begin{exercise}[$\star\star$]\label{exo:dv:build-a-pricing-library:4}
The standard deviation of the node \omterm{def:dv:greeks-and-the-hedging-pnl:greeks}{vega} is 51 with common random numbers. Where does that remaining noise come from, and how would you reduce it?
\end{exercise}

\begin{exercise}[$\star\star$]\label{exo:dv:build-a-pricing-library:5}
Why does the sum of the node \omterm{def:dv:greeks-and-the-hedging-pnl:greeks}{vegas} (6\,747) differ from the parallel \omterm{def:dv:greeks-and-the-hedging-pnl:greeks}{vega} (6\,706)?
\end{exercise}

\begin{exercise}[$\star\star$]\label{exo:dv:build-a-pricing-library:6}
What goes wrong if an engine reads a volatility from a global variable instead of from the snapshot?
\end{exercise}

\begin{exercise}[$\star\star\star$]\label{exo:dv:build-a-pricing-library:7}
\emph{Coding.} Register a Monte Carlo engine with 400\,000 paths and compute the node \omterm{def:dv:greeks-and-the-hedging-pnl:greeks}{vega} of the \omterm{def:dv:autocallables:autocallable}{autocallable} over twenty seeds with common random numbers. How does the
standard deviation scale?
\end{exercise}

\begin{exercise}[$\star\star\star$]\label{exo:dv:build-a-pricing-library:8}
\emph{Find the flaw.} ``Our new engine agrees with the old one to eight decimals on every trade, so it is correct.''
\end{exercise}

\section{Problem: The Vega Nobody Could Compute}

\begin{problem}[{Weekend problem --- one node, five pricers}]\label{pb:dv:build-a-pricing-library:1}
The risk manager of the opening wants the book's \omterm{def:dv:greeks-and-the-hedging-pnl:greeks}{vega} to the one-year 90\% node, today and every day.

\medskip\noindent\textbf{Part I --- The library.}
\begin{enumerate}
\item Which engine prices each of the five positions by default, and why that one?
\item What is the book's value?
\item How close are the engines on the one-year call, and what explains the gaps?
\item What must every engine do for the node bump to reach it?
\item What does the registry do when no engine supports a pair?
\end{enumerate}

\medskip\noindent\textbf{Part II --- The node.}
\begin{enumerate}[resume]
\item Give each position's \omterm{def:dv:greeks-and-the-hedging-pnl:greeks}{vega} to the node.
\item Why is the \omterm{def:dv:autocallables:autocallable}{autocallable}'s the largest, and why is its sign positive for the firm?
\item How does the put's strike of 95 give it half its exposure at 90?
\item What does the \omterm{def:dv:variance-swaps-and-volatility-derivatives:vs}{variance swap}'s strip contribute, and why at every node?
\item What is the book's total \omterm{def:dv:greeks-and-the-hedging-pnl:greeks}{vega} by expiry?
\end{enumerate}

\medskip\noindent\textbf{Part III --- The noise.}
\begin{enumerate}[resume]
\item Give the node \omterm{def:dv:greeks-and-the-hedging-pnl:greeks}{vega}'s standard deviation over twenty seeds with and without common random numbers.
\item Why does sharing random numbers cut it by a factor of 4.3?
\item What would the standard deviation be with four times the paths?
\item Why does the binding rule tie the seed to the instrument and not to the book?
\item What does the grid of ten scenarios say about the book?
\end{enumerate}

\medskip\noindent\textbf{Part IV --- Judgement.}
\begin{enumerate}[resume]
\item Would you report the node \omterm{def:dv:greeks-and-the-hedging-pnl:greeks}{vega} to one unit or to a hundred?
\item What would you add before Book 6's risk engine runs the library every night?
\item Which test would have caught a pricer that reads a stale surface?
\item State the \emph{named result}: the book's bucketed \omterm{def:dv:greeks-and-the-hedging-pnl:greeks}{vega} at the one-year 90\% strike, and the Monte Carlo noise in it with and without common random numbers.
\item In one sentence: what makes a Greek computable for every pricer?
\end{enumerate}
\end{problem}

\section{Interview questions}

\begin{interviewq}[$\star$ \iqroles{developer}]\label{iq:dv:build-a-pricing-library:1}
Why separate instruments, models and engines?
\end{interviewq}

\begin{interviewq}[$\star\star$ \iqroles{developer, risk}]\label{iq:dv:build-a-pricing-library:2}
Your Monte Carlo \omterm{def:dv:greeks-and-the-hedging-pnl:greeks}{Greeks} are noisy. What design rule fixes most of it?
\end{interviewq}

\begin{interviewq}[$\star\star$ \iqroles{developer}]\label{iq:dv:build-a-pricing-library:3}
Mutable market objects with observers, or immutable snapshots? Argue for one.
\end{interviewq}

\begin{interviewq}[$\star\star$ \iqroles{developer, researcher}]\label{iq:dv:build-a-pricing-library:4}
How do you test a new \omterm{def:dv:build-a-pricing-library:engine}{pricing engine}?
\end{interviewq}

\begin{interviewq}[$\star\star$ \iqroles{risk}]\label{iq:dv:build-a-pricing-library:5}
How would you compute bucketed \omterm{def:dv:greeks-and-the-hedging-pnl:greeks}{vega} for a book priced by different models and engines?
\end{interviewq}

\begin{interviewq}[$\star\star\star$ \iqroles{developer}]\label{iq:dv:build-a-pricing-library:6}
You must rewrite the core of the Python library in C++ without changing any number. How do you proceed?
\end{interviewq}
